%%%%%%%%%%  scan idx 84  |  folio 71  %%%%%%%%%%
\subsection*{Références bibliographiques~:}
\begin{enumerate}[label={[\arabic*]}]
\item \emph{Hirzebruch} : Neue topologischen Methoden in der algebraischen Geometrie.
\item \emph{Serre} : Faisceaux algébriques cohérents, Ann. of Maths.
\item \emph{Grothendieck} : Sur les faisceaux algébriques cohérents, in Séminaire Cartan 1956/57.
\end{enumerate}

\section*{Démonstration du théorème de Riemann--Roch}
\begin{center}
(Grothendieck, 1er Novembre 1957)
\end{center}

\subsection*{* Rappel de notations\footnotemark}
\footnotetext{Les passages entre astériques *\dots* sont des interpolations du copiste.}
[Ce sont celles du «~rapport sur Riemann--Roch~», cité R.R.R. dans ce qui suit.]

Toutes les variétés considérées sont non singulières et quasi-projectives (i.e.\ isomorphes à une sous-variété localement fermée d'un espace projectif). Le corps de base est algébriquement clos, de caractéristique quelconque.

Si \(X\) est une telle variété, on note \(K(X)\) le groupe abélien engendré par les classes de faisceaux cohérents sur \(X\) (modulo l'identification d'une extension à une somme)~; la somme alternée des Tor fait de \(K(X)\) un anneau. Si \(\underline F\) est un faisceau (resp.\ un fibré \(E\) à fibre vectorielle, resp.\ une sous-%
